\section{A finite regional replacement bound with moving modal patterns}
\label{sec:regional}

\paragraph{What is being closed.}
Equation~\eqref{eq:raybound} needed an unevaluated Taylor remainder.
Here the remainder is derived from the characteristic operator itself.
We bound sums of roots inside fixed contours, allowing their eigenvectors
and individual locations to change. No prescribed-pattern compatibility
condition is imposed on the regional upper bound. The construction uses
standard argument-principle, matrix-logarithm and convex-duality tools
\cite{beyn,lessard}; its specialization and usefulness for this grid
contract, rather than those tools, are the research question.

\subsection{A finite identity that survives internal root collisions}
\status{EXACT\_IDENTITY}
Use the full finite-dimensional characteristic matrix
$\Delta(s,p)$ of the linearized DDE, with fixed delays.
It is holomorphic in $s$ on and inside each positively oriented,
piecewise smooth simple contour $\Gamma$ used below.
Do not use a meromorphic Schur complement without accounting for its poles.
Exclude the rotational root from these contours. In the full state chart
this avoids requiring a parameter-independent gauge-deflation matrix.

Let $p=p_0+d$ lie in a declared compact design region $\mathcal D$ and
let $\Delta_0(s)=\Delta(s,p_0)$ be nonsingular on $\Gamma$.
Suppose a factorization
\[
 \Delta(s,p)-\Delta_0(s)=\mathcal U(s)\mathcal V(s,d),\qquad
 H(s,d)=\mathcal V(s,d)\Delta_0(s)^{-1}\mathcal U(s)
\]
is available with $H\in\mathbb C^{r\times r}$, $\mathcal V(s,0)=0$.
All dimensions and coordinates remain fixed on $\mathcal D$.
Assume the \emph{uniform} bound
\begin{equation}
 \sup_{d\in\mathcal D,\,s\in\Gamma}\|H(s,d)\|_2
 \le\kappa_\Gamma<1 . \label{eq:regional-kappa}
\end{equation}
The determinant lemma and invertibility of $I+tH$, $0\le t\le1$,
preserve the number $m_\Gamma>0$ of enclosed roots, counting multiplicity.
Define their mean $z_\Gamma(p)=m_\Gamma^{-1}
\sum_{\lambda\ {\rm inside}\ \Gamma}\lambda$.
Then
\begin{equation}
 \boxed{z_\Gamma(p)-z_\Gamma(p_0)
 =-\frac{1}{2\pi\ii m_\Gamma}
       \oint_\Gamma \operatorname{tr}\log(I+H(s,d))\,ds .}
 \label{eq:regional-log}
\end{equation}
\begin{proof}
The determinant ratio on the contour is $\det(I+H)$.
The norm-convergent series
$\operatorname{tr}\log(I+H)=
\sum_{\ell\ge1}(-1)^{\ell+1}\operatorname{tr}(H^\ell)/\ell$
chooses a continuous, single-valued logarithm along the closed contour.
The argument principle gives the difference of root sums as
$(2\pi\ii)^{-1}\oint s\,d\log\det(I+H)$.
Integration by parts gives \eqref{eq:regional-log}.
This only uses regularity on the contour and constant enclosed count;
roots may collide or become defective inside it.
\end{proof}

If every physical root satisfies $\Re\lambda\le-\sigma$, then necessarily
\[
 g_\Gamma(p):=\Re z_\Gamma(p)+\sigma\le0.
\]
The converse is false: a stable mean can hide an unstable member.
Several disjoint contours can sharpen the necessary conditions.
One-root contours recover individual sensitivities where regular.
A finite collection is sufficient for a \emph{necessary relaxation},
not for a full-spectrum feasibility claim.

\subsection{An explicit uniform remainder}
\status{PROVED\_REDUCED\_MODEL}
Write $H=H_1+H_{\rm rem}$ with $H_1$ linear in $d$, and suppose
\[
 \sup_{\Gamma,\mathcal D}\|H_{\rm rem}\|_*
 \le e_\Gamma,\qquad \|\cdot\|_* \text{ is the nuclear norm}.
\]
Let $\ell_\Gamma$ be the contour length and put
$\varphi(x)=-\log(1-x)-x$, $0\le x<1$.
The affine prediction and its remainder satisfy
\begin{align}
 g_\Gamma(p)&=g_\Gamma(p_0)+J_\Gamma d+\mathcal R_\Gamma(d),\\
 J_\Gamma d&=-\Re\frac{1}{2\pi\ii m_\Gamma}
                     \oint_\Gamma\operatorname{tr}H_1(s,d)\,ds,\\
 \boxed{|\mathcal R_\Gamma(d)|\le\beta_\Gamma
 :=\frac{\ell_\Gamma}{2\pi m_\Gamma}
       [e_\Gamma+r\varphi(\kappa_\Gamma)] .}
 \label{eq:regional-remainder}
\end{align}
\begin{proof}
The trace of $H-H_1$ is bounded by its nuclear norm.
Also $|\operatorname{tr}H^\ell|\le r\|H\|_2^\ell$,
including nonnormal $H$. Summing terms $\ell\ge2$ gives
$r\varphi(\kappa_\Gamma)$; contour integration proves the bound.
\end{proof}
A simpler conservative expression uses
$\varphi(x)\le x^2/[2(1-x)]$.
For the logarithm term, any consistently used induced matrix norm
can replace the Euclidean norm: each eigenvalue of $H$ is bounded
by that norm, hence $|\operatorname{tr}H^\ell|\le r\kappa^\ell$.
Pointwise verified bounds may instead be integrated along $\Gamma$.
Numerical errors in $g_\Gamma(p_0)$ and $J_\Gamma$ must be included:
if their certified errors are $e_0$ and $e_{J,a}$ and
$|d_a|\le r_a$, add $e_0+\sum_a e_{J,a}r_a$ to $\beta_\Gamma$.
No eigenvalue Hessian or fixed modal pattern is required.

\subsection{The actual replacement contract gives a small factorization}
\status{EXACT\_IDENTITY}
For the repository's exported fixed-support, fixed-equilibrium sharing
chart, define
\begin{align*}
 W_\rho&=\diag(1-\rho_1,1-\rho_1,\ldots,1-\rho_n,1-\rho_n),\\
 G(\rho)&=Y+W_\rho D_s,&
 N(\rho)&=W_\rho C_s+(I-W_\rho)C_f,\\
 V(\rho)&=-G(\rho)^{-1}N(\rho),&
 C(\rho)&=E_\theta+H_vV(\rho),\\
 B(K)&=B_p\diag K_p+B_i\diag K_i,&
 E_\tau(s)&=\diag(e^{-s\tau_i}).
\end{align*}
Here $V$ is the algebraic voltage \emph{Jacobian}, not the voltage phasor.
With ordinary gain coordinates,
\[
 \Delta(s,p)=sI-A_{\rm dev}-B_vV(\rho)-B(K)E_\tau(s)C(\rho).
\]
This is the inspected parametric model, not a universal identity for
all SG/GFL models or dispatch changes. If equilibrium/device matrices
change under another contract, their derivatives and remainders must
also be bounded. Endpoint state removal requires a separate chart.

Denote reference matrices by a subscript $0$, and set
\[
 J=G_0^{-1}\delta G,\quad
 V_1=-G_0^{-1}(\delta N+\delta G V_0),\quad
 \delta V=(I+J)^{-1}V_1,\quad
 V_{\rm rem}=- (I+J)^{-1}J V_1 .
\]
Thus $\delta V=V_1+V_{\rm rem}$ exactly. If uniformly
$\|J\|_2\le a<1$ and $\|V_1\|_2\le v$, then
\begin{equation}
 \|\delta V\|_2\le\frac{v}{1-a},\qquad
 \|V_{\rm rem}\|_2\le\frac{av}{1-a}. \label{eq:voltage-remainder}
\end{equation}
Because $G,N$ are affine in $\rho$, computable bounds are
$a=\sum_j r_{\rho,j}\|G_0^{-1}G_{\rho_j}\|_2$ and
$v=\sum_j r_{\rho,j}\|G_0^{-1}(N_{\rho_j}+G_{\rho_j}V_0)\|_2$.
These require verified solves and norm enclosures for a certificate.

Set $D_0(s)=B_v+B_0E_\tau(s)H_v$ and define
\begin{align}
 \mathcal U(s)&=[\,D_0(s)\quad B_pE_\tau(s)\quad B_iE_\tau(s)\,],\\
 \mathcal V(s,d)&=-
 \begin{bmatrix}
 \delta V\\
 \diag(\delta K_p)(C_0+H_v\delta V)\\
 \diag(\delta K_i)(C_0+H_v\delta V)
 \end{bmatrix}. \label{eq:physical-lowrank}
\end{align}
Direct expansion proves $\Delta-\Delta_0=\mathcal U\mathcal V$.
The linear part uses $V_1,C_0$ in the three block rows.
The remainder uses respectively
$V_{\rm rem}$, $\diag(\delta K_p)H_v\delta V$, and
$\diag(\delta K_i)H_v\delta V$, with the same minus signs.
For this export, $\mathcal U$ has 40 columns for 204 states:
20 voltage coordinates and two sets of ten PLL gain channels.
The dimension is not a discarded-state approximation.

For explicit conservative constants, let
$v_b=v/(1-a)$, $v_r=av/(1-a)$,
$k_p=\max_i|\delta K_{p,i}|_{\max}$,
$k_i=\max_i|\delta K_{i,i}|_{\max}$, and
$c_b=\|C_0\|_2+\|H_v\|_2v_b$.
Then
\begin{align*}
 \|\mathcal V\|_2&\le v_{\mathcal V}
   :=\sqrt{v_b^2+(k_pc_b)^2+(k_ic_b)^2},\\
 \|\mathcal V_{\rm rem}\|_2&\le v_{\rm rem}
   :=\sqrt{v_r^2+(k_p\|H_v\|_2v_b)^2+(k_i\|H_v\|_2v_b)^2}.
\end{align*}
If $w_\Gamma$ uniformly bounds
$\|\Delta_0(s)^{-1}\mathcal U(s)\|_2$ on the entire contour, use
$\kappa_\Gamma=v_{\mathcal V}w_\Gamma$ and
$e_\Gamma=r v_{\rm rem}w_\Gamma$.
Fixed nonsingular state/port scalings may improve these bounds if
applied consistently. A large bound is inconclusive, not instability.

\subsection{Stronger route: retain algebraic variables to remove that remainder}
\label{sec:affinelift}
\status{EXACT\_IDENTITY}
The inverse-network remainder can be avoided altogether for this
sharing contract. Retain the algebraic voltage perturbation $v$ and
the detector perturbation $e$ with the dynamic state $x$:
\begin{equation}
 \mathcal F(s,p)=
 \begin{bmatrix}
 sI-A_{\rm dev}&-B_v&-B(K)E_\tau(s)\\
 N(\rho)&G(\rho)&0\\
 -E_\theta&-H_v&I
 \end{bmatrix}. \label{eq:affineDAE}
\end{equation}
These are additional algebraic coordinates, not invented dynamics
or additional control channels. Eliminating $e$ and $v$ proves
\begin{equation}
 \det\mathcal F(s,p)=\det G(\rho)\det\Delta(s,p).
 \label{eq:descriptor-equivalence}
\end{equation}
On a region where $G$ is nonsingular, the finite roots and their
multiplicities coincide. The factor $\det G$ is independent of $s$.

Define $d=(\delta\rho,\delta K_p,\delta K_i)$ and the repeated diagonal
\[
 \Theta(d)=\diag(\delta\rho_1,\delta\rho_1,\ldots,
                \delta\rho_n,\delta\rho_n,\delta K_p,\delta K_i).
\]
With the reference operator $\mathcal F_0$, set
\begin{align}
 \mathcal U&=
 \begin{bmatrix}
 0&B_p&B_i\\ I&0&0\\0&0&0
 \end{bmatrix},&
 \mathcal V(s)&=-
 \begin{bmatrix}
 C_s-C_f&D_s&0\\
 0&0&E_\tau(s)\\
 0&0&E_\tau(s)
 \end{bmatrix},\\
 T(s)&=\mathcal V(s)\mathcal F_0(s)^{-1}\mathcal U,&
 H(s,d)&=\Theta(d)T(s). \label{eq:affine-return}
\end{align}
Then, exactly,
\[
 \mathcal F(s,p)-\mathcal F_0(s)
       =\mathcal U\Theta(d)\mathcal V(s),\qquad
 \frac{\det\mathcal F(s,p)}{\det\mathcal F_0(s)}
       =\det(I+\Theta(d)T(s)).
\]
The descriptor has dimension 234 and the determinant update has
dimension 40 for this export. In particular,
\[
 \boxed{H=H_1,\qquad H_{\rm rem}=0.}
\]
All nonlinear dependence of the finite root change is now in the
logarithm. The ordinary gain coordinates matter: log-gain increments
would require exponential parameter remainders.
Applying \eqref{eq:regional-log} to $\mathcal F$ gives the same
physical root means, and \eqref{eq:regional-remainder} simplifies to
\begin{equation}
 \boxed{\beta_\Gamma
   =\frac{\ell_\Gamma r}{2\pi m_\Gamma}
          \varphi(\kappa_\Gamma).} \label{eq:affine-remainder}
\end{equation}
The baseline-moment and coefficient integration errors still need
their separate enclosures.

\paragraph{Use the declared parameter box, not a product of large norms.}
Let $R_d$ contain the corresponding nonnegative half-widths, repeated
for the two occurrences of each replacement coordinate. Entrywise,
\[
 |H(s,d)|\le M(s):=R_d|T(s)|.
\]
For any fixed vector $w>0$ on a contour panel,
\begin{equation}
 \|H(s,d)\|_{w,\infty}
 \le \max_a\frac{(M(s)w)_a}{w_a}
 =:\kappa_w(s). \label{eq:positive-envelope}
\end{equation}
This follows by scaling with $\diag w$ and taking absolute row sums.
It is a standard positive-matrix bound, not a new stability theorem.
The repeated-coordinate correlation is relaxed conservatively;
the controllers remain local and the physical gain bounds unchanged.
A Perron-based weight can suggest $w$, but certification uses verified
inequalities for the chosen positive vector, not a floating-point
eigenvalue treated as an exact bound.

\paragraph{A concrete continuous-panel enclosure.}
On a panel contained in $|s-s_c|\le h$, put
$Q_c=\mathcal F_0(s_c)^{-1}$,
$W_c=Q_c\mathcal U$, and
\[
 e_\tau(h)=\max_i e^{-\Re(s_c)\tau_i}
                     (e^{h\tau_i}-1).
\]
For fixed delays,
\[
 \|\mathcal F_0(s)-\mathcal F_0(s_c)\|_2
 \le h+\|B_0\|_2e_\tau(h),\qquad
 \|\mathcal V(s)-\mathcal V(s_c)\|_2\le\sqrt2\,e_\tau(h).
\]
If verified bounds give
$a_c:=\|Q_c\|_2[h+\|B_0\|_2e_\tau(h)]<1$,
the Neumann series yields
\begin{equation}
 \|T(s)-T(s_c)\|_2
 \le
 \frac{\|\mathcal V(s_c)\|_2 a_c+\sqrt2\,e_\tau(h)}
      {1-a_c}\,\|W_c\|_2=:\epsilon_{T,c}. \label{eq:panel}
\end{equation}
Hence $R_d(|T(s_c)|+\epsilon_{T,c}\one\one^T)$ is an
entrywise panel envelope. Verified rounding/solve errors in
$Q_c,W_c,T(s_c)$ must be added; these formulas do not make an
ordinary dense solve self-certifying.
One positive-vector test on that envelope bounds the entire panel
and the entire parameter box. A finite covering closes the contour.
Thus the outstanding computation has an explicit error formula,
rather than an unspecified Hessian-sampling step.

\subsection{Regional maximum and a controller-eliminated witness}
\status{PROVED\_REDUCED\_MODEL}
Stack the necessary cluster inequalities and their verified errors.
Let $d=(\delta\rho,\delta K)$ range over compact boxes
$\mathcal R\times\mathcal K$, intersected with the original physical
limits. Let $A=J_\rho$, $C=J_K$, $b=-g_0+\beta$.
Every fully secure design in this region necessarily satisfies
\[
 A\delta\rho+C\delta K\le b.
\]
Consequently the LP
\begin{equation}
 U_{\mathcal D}=\max_{\delta\rho\in\mathcal R,\,
                         \delta K\in\mathcal K}
 (P^0)^T\delta\rho\quad
 \text{subject to }A\delta\rho+C\delta K\le b \label{eq:regionalLP}
\end{equation}
gives
\begin{equation}
 \boxed{P_F\le P^\star_{\mathcal D}
 \le (P^0)^T\rho_0+U_{\mathcal D}=P_{U,\mathcal D}.}
 \label{eq:regional-bracket}
\end{equation}
The lower candidate must belong to this same region and pass the full
spectrum and all declared events. Numerically validated feasibility
remains numerical evidence unless its own errors are certified.
Omitting frequency, RoCoF or actuator constraints from this LP is an
optimistic relaxation, not a safety claim.

For any $\eta\ge0$, the support functions of the two boxes give
\begin{equation}
 U_{\mathcal D}\le
 \eta^Tb+h_{\mathcal R}(P^0-A^T\eta)
          +h_{\mathcal K}(-C^T\eta). \label{eq:regionaldual}
\end{equation}
Minimization over $\eta$ recovers the feasible bounded LP value
by ordinary LP duality. For an interval box,
$h(z)=\sum_j\max(l_jz_j,h_jz_j)$.
The bound is valid for \emph{all} replacement directions in the
declared region, unlike the earlier fixed-ray calculation.

For a proposed replacement $\delta\rho$, the strict inequality
\begin{equation}
 \eta^T(A\delta\rho-b)>h_{\mathcal K}(-C^T\eta)
 \label{eq:regional-exclusion}
\end{equation}
proves that no permitted local gain vector in the region can satisfy
the necessary conditions, hence none can meet the full security contract.
The weights identify which collective requirements conflict.
If $C^T\eta=0$, their first-order aggregate has no gain authority;
finite residual authority is still allowed by $\eta^T\beta$.
Only the strict finite inequality, not a small derivative alone,
establishes exclusion. This is a regional statement, not a global
impossibility theorem for all local controllers.

\subsection{Where graph interactions actually enter}
\status{EXACT\_IDENTITY}
Group the rows/columns of $H$ by intervention site (two algebraic
voltage-equation channels and two gain channels per candidate).
The exact series contains
\[
 \operatorname{tr}H^2
 =\sum_i\operatorname{tr}H_{ii}^2
   +\sum_{i\ne j}\operatorname{tr}(H_{ij}H_{ji}),
\]
and $\operatorname{tr}H^\ell$ sums closed walks of length $\ell$
through these dynamical intervention ports. Their weights include
the full resolvent, P/Q coupling, SG dynamics and exact delay phases.
They are not necessarily physical transmission-line cycles, and $H$
is neither a Laplacian nor a physical damping matrix.
Even its diagonal blocks depend on the full network.

Writing $H=D+O$ with $D$ block diagonal gives
\[
 \det(I+H)=\det(I+D)\det(I+T),\qquad T=(I+D)^{-1}O,\quad
 \operatorname{tr}T=0.
\]
When the relevant logarithm series converge, the first omitted
inter-site term after discarding $O$ is
$-\tfrac12\operatorname{tr}T^2$ in the log determinant.
Its contribution to the root-sum change has the opposite sign
because of \eqref{eq:regional-log}; it need not be stabilizing.
If $\|T\|_2\le t<1$, the discarded log-determinant magnitude is
at most $r\varphi(t)$.
This supplies a falsifiable approximation error for neglecting
collective interactions, rather than assuming graph roughness
predicts damping.

An orthogonal graph basis can reorganize these operators.
It preserves traces and Euclidean norms for the corresponding
similarity, but its diagonal truncation changes the approximation.
Low graph frequency alone does not justify that truncation.
No GSP advantage or physical cycle mechanism is established until
the retained/off-diagonal terms explain a measured limiting decision.

\subsection{What must be verified before any IEEE--39 certificate}
The closed formulas above replace an unspecified eigenvalue Hessian
by explicit matrix-inverse and logarithm remainders. The affine
descriptor removes the matrix-parameter remainder entirely for
the inspected sharing contract. These identities do not
automatically make a floating-point evaluation rigorous.
For each regional box, an executable verification must:
\begin{enumerate}
 \item Verify the equilibrium/parameter chart and algebraic inverse
 over that box, including any input/linearization uncertainty.
 \item Enclose $\Delta_0^{-1}\mathcal U$ on every \emph{continuous}
 contour panel with outward-rounded complex arithmetic and verified
 linear solves; prove the baseline enclosed count and boundary regularity.
 \item Enclose the baseline moments, coefficient integrals and
 $\kappa_\Gamma<1$ uniformly. A sampled maximum does not do this.
 \item Form the optimistic LP, bound numerical coefficient/dual errors,
 and compare its upper bound with a fully checked candidate in the box.
 \item Subdivide only to tighten the declared region; a global bound
 needs exhaustive coverage of the original domain, including chart changes.
\end{enumerate}
If a condition fails, report an unresolved box or a conservative bound.
The average-root condition can be weak even with a perfect remainder.
A candidate with all strict security slacks and interior replacement
fractions is not an established local maximum: continuity gives a
small further feasible replacement if those strict inequalities hold
for the exact model. The previously validated point therefore does
not itself establish the presence of a modal conflict.

\paragraph{Current status and novelty boundary.}
The identities and conditional regional relaxation have proofs.
Small symbolic and algebraic checks are recorded separately in
\path{regional_certificate_20261003/}.
No uniform IEEE--39 contour/parameter enclosure, nontrivial MW upper
bound, controller-exclusion witness, or advantage over conventional
tuning has been obtained in this extension.
Contour eigensolvers, validated DDE spectra, and direct delay-system
controller optimization already exist \cite{beyn,lessard,michielsdesign}.
The candidate contribution is an informative, physically interpreted
replacement bracket under the actual PLL-only action set. Its novelty
and usefulness remain to be demonstrated; the proofs alone do not
establish either.
